%% Shared problem body, included by both `problem.rbx.tex` (the standalone,
%% full-document template) and `problem-fragment.rbx.tex` (the fragment that is
%% `\subimport`-ed into the contest task sheet). Keeping it in one place means
%% the standalone PDF and the joined contest PDF always render a problem the
%% same way. Edit the look of every problem here.
\includeProblem
%- if problem.short_name is defined
	[\VAR{problem.short_name}]
%- endif
{\VAR{problem.title | escape}}{

%- if params.show_limits is truthy
	\includeLimits
		{\VAR{problem.limits.timeLimit} ms}
		{\VAR{problem.limits.memoryLimit} MiB}
		{}
%- else
	\vspace{0.55cm}
%- endif

	%- if problem.blocks.macros is defined
		\VAR{problem.blocks.macros}
	%- endif

	\VAR{problem.blocks.legend}

	%- if problem.blocks.input is defined
		\inputdesc{\VAR{problem.blocks.input}}
	%- endif

	%- if problem.blocks.output is defined
		\outputdesc{\VAR{problem.blocks.output}}
	%- endif

	%- if problem.blocks.interaction is defined
		\interactiondesc{\VAR{problem.blocks.interaction}}
	%- endif

	%- if problem.samples
		\vspace{0.1cm}
		\subsection*{\strExamples}
		%- for sample in problem.samples
			%- if sample.interaction is not none
				\exampleInteractive
				%- for chunk in sample.interaction.chunks
					\interactionchunk{\VAR{chunk.path}}{\VAR{chunk.pipe}}
				%- endfor
			%- else
				%% `sample.input`/`sample.output` are root-relative paths to the
				%% generated I/O (read verbatim by `\VerbatimInput` inside `\example`).
				\example{\VAR{sample.input}}
					{\VAR{sample.output if sample.output is not none else ''}}
			%- endif
			%- if sample.explanation_file is not none
				%% The explanation lives in its own sample folder and is pulled in
				%% with `\subimport` so it can carry its own figures.
				\explanation{\subimport{\VAR{sample.dir}}{\VAR{sample.explanation_file}}}
			%- endif
		%- endfor
	%- endif

	%- if problem.blocks.notes is defined
		\subsection*{\strNotes}
		\VAR{problem.blocks.notes}
	%- endif
}
